\documentclass[11pt]{article}
\usepackage[margin=1in]{geometry}
\usepackage{titlesec}
\usepackage{enumitem}
\usepackage{xcolor}
\usepackage{tabularx}
\usepackage{booktabs}

\definecolor{umblue}{RGB}{0,48,86} % University of Mannheim blue
\titleformat{\section}{\Large\bfseries\color{umblue}}{\thesection}{1em}{}
\setlist{noitemsep,topsep=2pt}

\begin{document}

\section*{Ehrenwörtliche Erklärung}

Ich versichere, dass ich die beiliegende Bachelor-, Master-, Seminar-, oder
Projektarbeit ohne Hilfe Dritter und ohne Benutzung anderer als der angegebenen
Quellen und in der untenstehenden Tabelle angegebenen Hilfsmittel angefertigt
und die den benutzten Quellen wörtlich oder inhaltlich entnommenen Stellen als
solche kenntlich gemacht habe. Diese Arbeit hat in gleicher oder ähnlicher Form
noch keiner Prüfungsbehörde vorgelegen. Ich bin mir bewusst, dass eine falsche
Erklärung rechtliche Folgen haben wird.

% Declare below which AI tools you used in the process of writing your work,
% including text, image, code, and data generation.
%
% The AI declaration must include at least one an entry for every graded subtask
% (i.e., where points are specified). The entry should state whether or not you
% used AI to solve parts of the task. Using AI for things such as language
% improvements or API search does not count as "Solved with AI", but things such
% as code gen, writing drafts, or solving math tasks does.
%
% Tasks that you (partially or fully) solved with AI will be graded more
% stringently. So make sure that whenever you use AI, the results are correct.
% My recommendation is to not use AI to solve the tasks, as this undermines the
% learning goals of this course and, ultimately, your own education.

\begin{center}
  \textbf{Declaration of Used AI Tools and LLMs} \\[.3em]
  \begin{tabularx}{\textwidth}{lclXc}
    \toprule
    Task & Solved with AI? & Tool & Purpose & Useful? \\
    \midrule
    1.1.1 & Yes & Claude Code &  Understanding parts of equations, (e.g. chain rule/ partial derivative of beta). & yes \\
    1.1.2 & Yes & Claude Code & Understanding parts of equations, (e.g. chain rule/ partial derivative of beta). Giving hints into the right direction when developing proofs after solving ß0, correcting proactive ideas of the author. & yes \\
    1.2.1 & No & & & \\
    1.2.2 & No & Internet websites & & yes\\
    1.2.3 & No & (see Jupyter Notebook) & & \\
    2.1 & No & & & \\
    2.2 & No & & & \\
    2.3 & No & & & \\
    \bottomrule
  \end{tabularx}
\end{center}

\vspace{2cm}
\noindent Unterschrift\\

\noindent Mannheim, den 12.09.2026 \hfill

\end{document}
